% Generated draft: requires mathematical and source review.
\documentclass[11pt]{article}
\usepackage[margin=25mm]{geometry}
\usepackage{fontspec}
\setmainfont{Latin Modern Roman}
\usepackage{amsmath,amssymb,mathtools,unicode-math}
\setmathfont{Latin Modern Math}
\usepackage{enumitem,needspace,xurl,hyperref}
\setlength{\parindent}{0pt}
\setlength{\parskip}{5pt}
\setlength{\emergencystretch}{3em}
\newcommand{\N}{\mathbb{N}}
\newcommand{\Z}{\mathbb{Z}}
\newcommand{\Q}{\mathbb{Q}}
\newcommand{\R}{\mathbb{R}}
\newcommand{\C}{\mathbb{C}}
\newcommand{\F}{\mathbb{F}}
\newcommand{\A}{\mathbb{A}}
\newcommand{\PP}{\mathbb{P}}
\newcommand{\E}{\mathbb{E}}
\newcommand{\eps}{\varepsilon}
\newcommand{\dd}{\,\mathrm{d}}
\begin{document}
\section*{UnsolvedMath ID: 11000057}
\textbf{Source problem code:} catalogue entry
\subsection{Definitions and mathematical statement}
Let \(I_g\) be the Torelli group, in the usual closed oriented genus-\(g\) convention the kernel of \(\operatorname{Mod}_g\to\operatorname{Sp}(H_1(\Sigma_g;\Z))\). A \(K(I_g,1)\) is a connected CW complex with fundamental group \(I_g\) and contractible universal cover. For an integer \(n\geq0\), say \(I_g\) has type \(F_n\) when such a complex has finitely many cells in dimensions at most \(n\). Define
\[
 f(g)=\sup\{n\geq0:I_g\text{ has type }F_n\}\in\Z_{\geq0}\cup\{\infty\}.
\]
The problem is to determine this finiteness length (for the genera in the source, in particular \(g\geq2\)). Type \(F_1\) means finite generation and type \(F_2\) finite presentation; the question includes all higher skeleton finiteness conditions, not just those two.

The source reports:
\begin{enumerate}
\item \(f(2)=0\), from failure of finite generation, attributed to McCullough--Miller [McM].
\item \(f(3)\leq3\), attributed to unpublished Johnson--Millson work referred to in [Me].
\item For \(g\geq3\), \(1\leq f(g)\leq6g-5\), combining Johnson's finite-generation result [Jo2] with Akita's theorem, numbered 5.13 later in the source.
\end{enumerate}
It proposes the guess \(f(g)=g-2\), qualified as fitting the limited known data. All reported bounds are dated source statements, not independently updated claims. A following folklore special-case conjecture is announced but not supplied.
\subsection{Short English statement}
Determine the Torelli group's finiteness length: the largest skeleton dimension up to which a classifying space can have finitely many cells, with the source suggesting \(f(g)=g-2\).
\subsection{Sources}
\begin{itemize}
\item \emph{Problems on Mapping Class Groups and Related Topics}, edited by Farb, Problem 5.11, ``Torelli finiteness,'' supplied PDF page 47. \url{https://www.math.uchicago.edu/~farb/papers/mcgbook.pdf}. Recorded access 2026-07-29; extraction pdf-font-structure.
\item McCullough--Miller [McM]; Johnson--Millson, unpublished, as referred to in [Me]; Johnson [Jo2]; and Akita, source Theorem 5.13. Complete bibliography entries and the latter theorem statement are absent.
\item Local UnsolvedMath registry, \texttt{data/unsolvedmath/problems.json}, permanent ID \(11000057\), exact source code \texttt{catalogue entry}. The OPEN-TRIAGE summary also cross-refers to \texttt{AMR-109-0002}, \texttt{AMR-109-0058} and \texttt{AMR-109-0305}, without supplying their statements.
\end{itemize}
\end{document}
